\section{Related Work}
In this section, we review the related work in three aspects, namely Reciprocal Recommendation, Sequential Recommendation and Self-Distillation.

\subsection{Reciprocal Recommendation}
Reciprocal recommender system (RRS)~\cite{pizzato2010recon} is a recommender system between people and people that significantly deviates from conventional user-item recommender systems. 
To emphasize the mutual matching relationship between users, RRS requires modeling the interests of both parties in the recommendation, which poses a greater challenge for delivering effective recommendations.
Existing RRS studies can be roughly categorized into two groups: content-based methods~\cite{pizzato2010recon,alanazi2013people,tu2014online} and collaborative filtering-based methods~\cite{cai2012reciprocal,xia2015reciprocal}. Content-based methods focus on the matching of user profiles,
while collaborative filtering-based methods pay more attention to modeling historical user interactions. 
There also exist hybrid models that integrate user behavior modeling as well as content matching~\cite{akehurst2011ccr,yang2017combining,kleinerman2018optimally,hou2022leveraging}. 
In recent years, deep learning techniques have been increasingly used to improve the expressive power of RRS models, including matrix factorization~\cite{qu2018reciprocal}, latent factor models~\cite{neve2019latent}, factorization machines~\cite{yildirim2021bideepfm}, and graph neural networks (GNNs)~\cite{luo2020motif,yang2022modeling} 
Furthermore, several studies also make use of multiple types of user behaviors to learn better bilateral user preferences in RRS scenarios~\cite{le2019towards,fu2021beyond}. 
Although effective to some extent,
existing methods mainly encode user behaviors and profiles into static representations to model long-term user interests, ignoring the dynamic changes over time. In this work, we formulate reciprocal recommendation as a distinctive sequence matching task and propose to leverage bilateral behavior sequences for dynamic interest modeling on both sides. 


\subsection{Sequential Recommendation} \label{sec:seqrec}
In contrast to conventional collaborative filtering algorithms based on all historical user-item interactions~\cite{koren2008factorization,kabbur2013fism,he2018nais,rendle2009bpr,he2017neural,he2020lightgcn}, sequential recommendation pays more attention to the temporal aspects of user behaviors~\cite{hidasi2015session,sun2019bert4rec,kang2018self}.
Early sequential recommendation methods are frequently founded upon Markov Chains techniques~\cite{rendle2010factorizing,he2016fusing}. With the fast development of deep neural networks, numerous deep models like Recurrent Neural Networks (RNNs) are employed to facilitate item sequence modeling~\cite{hidasi2015session,tang2018personalized,yuan2019simple}. 
More recently, GNN-based~\cite{wu2019session,xu2019graph} and Transformer-based methods~\cite{kang2018self,sun2019bert4rec,hou2022core} also show great effects on sequential recommendation task. 
Besides direct item sequence modeling, recent studies also exploited techniques to additionally enhance model effectiveness and robustness, such as pre-training~\cite{zhou2020s3,hou2022towards}, data augmentation~\cite{wang2021counterfactual} and regularized training~\cite{xie2022cl4rec,zhou2023mvs}.

The above sequential recommendation methods
have achieved amazing performance in conventional user-item recommendation. 
Nevertheless, in the reciprocal recommendation scenario, only single-perspective modeling usually fails to achieve satisfactory results. 
The other possible reason is that compared with the static item target in user-item recommendation, the dynamic characteristics of both users in reciprocal recommendation require more detailed interactions between both parties.
In our approach, we model user behavior sequences from both active and passive perspectives, and a multi-scale matching prediction technique is proposed for comprehensive and fine-grained two-way interactions.

\subsection{Self-Distillation}
Knowledge distillation~\cite{hinton2015distilling} is a popular technique that can effectively transfer knowledge from complex teacher networks to more streamlined student networks. In the realm of recommender systems, it is often used for domain knowledge transfer~\cite{yang2022cross,tu2021conditional} or model compression~\cite{kang2021topology,tian2023directed}. However, such methods mainly necessitate meticulously crafted teacher models, which incur additional design and experimental costs. 
In recent years, the self-distillation technique~\cite{zhang2019your} that distills knowledge from itself has been proposed, and there are three prevalent self-distillation methods exist: \textbf{(a)} distillation from deep to shallow layers of the same network~\cite{zhang2019your,luan2019msd,liu2020fastbert}; \textbf{(b)} using the self-ensemble model (from the most recent time steps) as a virtual teacher model~\cite{xu2020improving,shen2022self}; \textbf{(c)} data augmentation-based methods~\cite{huang2020comprehensive,xu2019data}.
In comparison to conventional knowledge distillation, self-distillation remains an area that is worth exploring within the realm of recommendation systems~\cite{xia2022device}.
In order to ensure low time consumption while engaging in more fine-grained two-way interactions, we put forth a simple self-distillation framework in this paper.
Our approach is most similar to the first of the above self-distillation methods, but instead of using the deeper network outputs as knowledge, we use time-sensitive fine-grained micro-level matching as a teacher for efficient macro-level matching.






